\begin{claim}{D.40}{Proved}{A finite homomorphism-count obstruction to the extra deck symmetry}
The mixed-multiplier four-cusp involution constructed in D.29
commutes with no single-transposition involution of multiplier $1$ in its
full loop-outer group. In particular it admits none of the
specified geometric deck actions required for the descent of D.37,
including actions moving the two-pants separator.

More precisely, for the exact profinite group
\[
 E=\overline{\langle a,l,t\mid
       tlt^{-1}=l^\epsilon,\ t^2=l^\kappa\rangle},
 \qquad p=a^{-1}l,
                                                        \tag{D.40a}
\]
with $\epsilon,\kappa$ as in D.29 and $\lambda=-\epsilon$,
there is no continuous automorphism sending the conjugacy
classes of $a,p$ to those of $a^\lambda,p$, respectively.
There is also none sending those classes to $a,p^\lambda$.
The obstruction is the exact inequality
\[
 \begin{aligned}
 \#\{\phi:E\to Q:\phi(a),\phi(p)\in C\}&=10200,\\
 \#\{\phi:E\to Q:\phi(a)\in C^{-1},\phi(p)\in C\}&=4920,
 \end{aligned}
                                                        \tag{D.40b}
\]
where all homomorphisms are continuous,
$Q=\operatorname{PSL}_2(\F_{11})$, and $C$ is the
conjugacy class of
$\left(\begin{smallmatrix}0&1\\10&2\end{smallmatrix}\right)$
modulo scalar matrices $\{\pm I\}$.
\end{claim}
\begin{proof}
We first identify the group and its peripheral marking.
In the basis $(l,a,b')$ of (D.29.10), adjoining the outer
lift $t$ gives the relations
\[
 tlt^{-1}=l^\epsilon,\quad tat^{-1}=b',\quad
 tb't^{-1}=l^\kappa a l^{-\kappa},\quad t^2=l^\kappa.
\]
Eliminating $b'$ gives (D.40a): the third relation follows
from the second and fourth. Conversely $b'=tat^{-1}$
defines the inverse homomorphism from (D.40a) to this
presentation. Both composites fix the named generators.
Thus this is an isomorphism of profinite groups, and
\[
                         E=N_{\rm bad}\amalg\widehat\Z(a)
                                                        \tag{D.40c}
\]
with the first factor generated by $l,t$.
The index-two subgroup $H$ is the kernel of the character
sending $t$ to $1$ and $a,l$ to $0$ in $C_2$.

The four $H$-peripheral classes of D.29 are represented by
\[
 a,\quad p=a^{-1}l,\quad b,\quad b^{-1}l^{-1}.
\]
Since $\alpha$ interchanges the first and third classes,
and the second and fourth, with multiplier $\lambda$,
and conjugation by $t$ is $\alpha$, these are equally
the classes of
\[
 a,\quad p,\quad (tat^{-1})^\lambda,
                         \quad(tpt^{-1})^\lambda.       \tag{D.40d}
\]
Here $\lambda^2=1$, and the original cusp conjugators
belong to $H$; thus (D.40d) records actual conjugacy
classes under the specified inclusion into $E$.
This description is independent of the Galois conjugator
$f$ used to construct the two pants isomorphisms.

For a finite group $Q$ define the central function
\[
 F_Q(L)=\#\{T\in Q:T L T^{-1}=L^\epsilon,
                                      \ T^2=L^\kappa\}.
                                                        \tag{D.40e}
\]
Every assignment $a\mapsto A$, $p\mapsto P$ forces
$l\mapsto AP$. The presentation (D.40a) then gives
exactly $F_Q(AP)$ continuous homomorphisms to $Q$.
For two conjugacy classes $C,D$ put
\[
             K_Q(C,D)=\sum_{A\in C,\ P\in D}F_Q(AP).
                                                        \tag{D.40f}
\]
Conjugating $L$ conjugates all its possible $T$, so
$F_Q$ is central. Since $PA=A^{-1}(AP)A$, this also proves
$K_Q(C,D)=K_Q(D,C)$.

Precomposition by a continuous automorphism of $E$ is
a bijection on its continuous homomorphisms to $Q$.
An automorphism with the first pair of class images in
the claim would therefore give
$K_Q(C,D)=K_Q(C^\lambda,D)$ for every pair $C,D$.
The inverse map has the same power effect because
$\lambda^2=1$, so this is equality of the two full sets,
without a choice of conjugators in the target.
The second proposed class effect would similarly give
$K_Q(C,D)=K_Q(C,D^\lambda)$.

We compute these integers in one completely specified
finite group. Enumerate the determinant-one matrices
over $\F_{11}$, and identify each matrix with its
negative, using their lexicographically smaller
representative. Matrix multiplication followed by this
identification constructs all $660$ elements of $Q$.
The class $C$ has size $60$ and order $11$.
Its inverse class is distinct and has the same size.
Since $\lambda_{11}=-1$, we have $C^\lambda=C^{-1}$.
The following table gives the full product counts needed
in (D.40f). The representative of each nonidentity
column is $\left(\begin{smallmatrix}0&1\\10&r\end{smallmatrix}\right)$;
the last column is the identity. We write
$N_{C,D}(L)=\#\{(A,P)\in C\times D:AP\in[L]\}$.
\[
\begin{array}{c|rrrrrrrr}
 r&0&1&2&3&4&5&9&I\\\hline
 |L|&2&3&11&5&5&6&11&1\\
 \text{class size}&55&110&60&132&132&110&60&1\\
 F_Q(L)&0&6&1&1&1&0&1&56\\
 N_{C,C}(L)
   &0&1320&120&660&660&0&840&0\\
 N_{C^{-1},C}(L)
   &660&0&120&660&660&1320&120&60
\end{array}                                               \tag{D.40g}
\]
Multiplying either of the last two rows by the
$F_Q$ row and summing gives respectively
\[
 7920+120+660+660+840=10200,
 \qquad 120+660+660+120+3360=4920.
\]
These unequal integers prove the first nonexistence
assertion. Symmetry of $K_Q$ proves the second.

All entries of (D.40g) are verified by the finite
enumeration in \path{code/root/haran_pair_counts.py}.
It constructs the multiplication, inverse, power and
conjugacy tables directly from the displayed matrices,
computes the residues of $\epsilon,\kappa$ modulo each
element order, and tests both equations in (D.40e) for
every pair $(L,T)$. It then counts each product in the
two displayed class pairs. Two independent computations,
\path{code/part_a/haran_pair_count_audit.py} and
\path{code/part_c/deck_obstruction_psl11_independent.py},
instead generate the faithful twelve-point permutation
action by $z\mapsto z+1$ and $z\mapsto-1/z$ on
$\mathbf P^1(\F_{11})$. The first enumerates allowed
$(L,T)$ before setting $P=A^{-1}L$; the second also
enumerates the triples directly. Both reproduce (D.40b).
The attached outputs include the full integer data.

Finally suppose a positive single-transposition
involution $\delta$ commuted with the D.29 involution
$\tau$. Their distinct peripheral permutations make
$\langle\delta,\tau\rangle\simeq V_4$.
Its center-free pullback extension inside $\Aut(H)$
contains $E$ as the normal preimage of $\langle\tau\rangle$,
using the canonical conjugation identification of
center-free extensions. Conjugation by any lift of
$\delta$ therefore defines a continuous automorphism
$\eta$ of this exact $E$.
By D.37's permutation computation, $\delta$ is either
$(1\ 3)$ or $(2\ 4)$ in the labeling (D.40d).
Its positive multiplier is $1$.
In the first case $\eta$ sends $a$ to an $E$-conjugate
of $a^\lambda$ and $p$ to an $E$-conjugate of $p$;
in the second case the roles are reversed.
Both have just been excluded. The argument imposes
no preservation of the original pants separator and
no order condition on the chosen lift of $\delta$.
It proves the stated obstruction for this actual
four-cusp construction; the general negative
three-cusp conjugacy question remains in D.12.
\end{proof}
